\documentclass[11pt]{article}

\usepackage[T1]{fontenc}
\usepackage{lmodern}
\usepackage{amsmath,amssymb,amsthm}
\usepackage{booktabs}
\usepackage{microtype}
\usepackage[hidelinks]{hyperref}
\usepackage[margin=1in]{geometry}

\newtheorem{theorem}{Theorem}
\newtheorem{lemma}[theorem]{Lemma}
\newtheorem{proposition}[theorem]{Proposition}
\theoremstyle{definition}
\newtheorem{definition}[theorem]{Definition}

\newcommand{\E}{\mathbb{E}}
\newcommand{\cE}{\mathcal{E}}
\newcommand{\gap}{\operatorname{gap}}

\title{An Exact Computational Lower Bound $m_2(4)\geq 18$\\
for a Width-Two Chain-Gap Threshold}
\author{Xiaojian Zeng\\
\small PBC School of Finance, Tsinghua University, Beijing 100190, China\\
\small \texttt{xiaojianzeng@foxmail.com}}
\date{7 September 2026}

\begin{document}
\maketitle

\begin{abstract}
For a finite poset $P$, let $h_P(x)$ be the expected rank of $x$ in a
uniformly random linear extension.  The chain gap of an inclusion-maximal
chain is the largest spacing in its expected ranks after adjoining the
endpoints $0$ and $|P|+1$.  Let $m_2(4)$ be the least order of a poset of
width exactly two for which every nonempty inclusion-maximal chain has chain
gap strictly greater than $4$.  We prove the exact computational lower bound
$m_2(4)\geq18$.  The proof exhausts all two-chain presentations through order
$17$, tests expected ranks using integer path counts, and reduces the
universal maximal-chain condition to exact reachability in the cover graph.
Altogether, $346{,}767{,}961$ width-two presentations are tested, with no
qualifier.  A structurally different verifier exhaustively agrees through
order $6$.  This is a partial result: order $18$ is not tested, the value of
$m_2(4)$ is not determined, and no minimizers are classified.
\end{abstract}

\section{The finite-threshold question}

Let $P=(V,\leq_P)$ be a finite nonempty poset of order $n$.  A linear
extension is an order-preserving bijection $f:V\to\{1,\ldots,n\}$.  With the
uniform distribution on the finite set $\cE(P)$ of linear extensions, put
\[
  h_P(x)=\E[f(x)]
  =\frac{1}{|\cE(P)|}\sum_{f\in\cE(P)}f(x).
\]
If $C=\{y_1<_P\cdots<_P y_m\}$ is a nonempty chain, define
\[
 \gap_P(C)=\max\{h_P(y_1),h_P(y_2)-h_P(y_1),\ldots,
 h_P(y_m)-h_P(y_{m-1}),n+1-h_P(y_m)\}.
\]
This is the chain-gap definition of Haqi~\cite{haqi2026gap}.  Here and below,
``maximal chain'' means inclusion-maximal, not maximum-cardinality.

Write $w(P)$ for the width of $P$ and set
\begin{align*}
 m_2(4)=\min\bigl\{n:{}&\text{there is an order-$n$ poset $P$ with }w(P)=2,\\
 &\gap_P(C)>4\text{ for every maximal chain }C\bigr\}.
\end{align*}
Haqi's construction gives, for every $L>0$, a width-two poset for which every
nonempty chain has gap at least $L$~\cite[Theorem~5.1]{haqi2026gap}; hence the
minimum exists.  That asymptotic existence theorem does not locate the first
order.  The contribution here is the following finite exclusion.

\begin{theorem}\label{thm:main}
No poset of width exactly two and order at most $17$ has
$\gap_P(C)>4$ for every nonempty inclusion-maximal chain $C$.  Consequently,
\[
  m_2(4)\geq18.
\]
\end{theorem}

The theorem does not decide whether an order-$18$ example exists and does not
classify minimizers.  A bounded search of the nearby literature located the
existence theorem above but no computation of this finite threshold; this is
not a claim of priority.

The computational scope is deliberately distinguished from a census of all
unlabeled posets.  It is enough for nonexistence to enumerate a redundant
family containing every width-two isomorphism class.  Related exact finite
work has instead combined symmetry-reduced enumeration with exact predicates
and an independently reconstructing verifier~\cite{zeng2026finite}; that
certificate architecture motivates the separate small-order check described
in Section~\ref{sec:verification}, but no result from that paper is used in
the proof.

\section{Complete two-chain encoding}

Fix positive integers $a\leq b$ with $a+b=n$, and consider two displayed
chains
\[
 A_0<\cdots<A_{a-1},\qquad B_0<\cdots<B_{b-1}.
\]
For each $i$, encode all cross-relations by
\[
 L_i=|\{j:B_j<A_i\}|,
 \qquad
 U_i=\min\bigl(\{j:A_i<B_j\}\cup\{b\}\bigr).
\]

\begin{lemma}[Two-chain boundary encoding]\label{lem:encoding}
The two-chain presentations are in bijection with pairs of sequences satisfying
\[
 0\leq L_0\leq\cdots\leq L_{a-1}\leq b,
 \qquad
 0\leq U_0\leq\cdots\leq U_{a-1}\leq b,
 \qquad L_i\leq U_i.
\]
For such a pair, the cross-relations are precisely
\[
 B_j<A_i\iff j<L_i,\qquad A_i<B_j\iff j\geq U_i.
\]
The presented order has width exactly two if and only if $L_i<U_i$ for at
least one $i$.
\end{lemma}

\begin{proof}
In a two-chain presentation, transitivity along the $A$-chain makes both
boundary sequences nondecreasing, and antisymmetry gives $L_i\leq U_i$.
Conversely, the displayed rules and monotonicity make every relation forced
by either chain transitive across the other chain; $L_i\leq U_i$ prevents a
cross-cycle.  Thus they define a partial order and recover the stated
boundaries.

If all $L_i=U_i$, every $A_i$ is comparable with every $B_j$, so the order is
a chain.  If some $L_i<U_i$, the pair $A_i,B_{L_i}$ is incomparable; the
displayed two-chain cover shows width at most two, so the width is exactly
two.
\end{proof}

By Dilworth's decomposition theorem~\cite{dilworth1950}, every finite
width-at-most-two poset is the union of two chains.  For width exactly two
these can be taken nonempty, and swapping them gives $a\leq b$.  Therefore,
as $a$ ranges from $1$ to $\lfloor n/2\rfloor$, Lemma~\ref{lem:encoding}
generates every width-two isomorphism class, possibly more than once.  Such
duplication cannot invalidate a nonexistence result.  It would, however, be
insufficient for the minimizer-classification part of the original question.

\section{Exact ranks from lattice paths}

For a fixed presentation, a state $(i,j)$ records an extension prefix
containing $A_0,\ldots,A_{i-1}$ and $B_0,\ldots,B_{j-1}$.  Let
\[
 K_j=|\{i:U_i\leq j\}|.
\]
From $(i,j)$ the next $A$-element is available exactly when $i<a$ and
$j\geq L_i$; the next $B$-element is available exactly when $j<b$ and
$i\geq K_j$.  Thus linear extensions are precisely the feasible monotone
lattice paths from $(0,0)$ to $(a,b)$.

Let $F(i,j)$ be the number of feasible prefixes ending at $(i,j)$ and
$G(i,j)$ the number of feasible suffixes beginning there.  They are computed
by the preceding two transition rules, with $F(0,0)=G(a,b)=1$.  In
particular,
\[
 e(P):=|\cE(P)|=F(a,b).
\]
If $S(x)=\sum_{f\in\cE(P)}f(x)$ is the integer rank sum, splitting paths at
the unique step that inserts an element gives
\begin{align}
 S(A_i)&=\sum_{j=L_i}^{b}F(i,j)G(i+1,j)(i+j+1),\label{eq:rankA}\\
 S(B_j)&=\sum_{i=K_j}^{a}F(i,j)G(i,j+1)(i+j+1).\label{eq:rankB}
\end{align}
Hence $h_P(x)=S(x)/e(P)$ is obtained without floating-point arithmetic.

\begin{lemma}[Exact maximal-chain test]\label{lem:chain}
A finite poset $P$ has a maximal chain $C$ with $\gap_P(C)\leq4$ if and only
if its cover graph has a directed path $x_1\lessdot\cdots\lessdot x_t$ from
a minimal element to a maximal element such that
\[
 S(x_1)\leq4e(P),\quad
 S(x_{r+1})-S(x_r)\leq4e(P),\quad
 (n+1)e(P)-S(x_t)\leq4e(P).
\]
\end{lemma}

\begin{proof}
An inclusion-maximal chain in a finite poset begins at a minimal element,
ends at a maximal element, and uses only cover relations; conversely every
such saturated path is an inclusion-maximal chain.  Multiplying all expected
ranks by the positive integer $e(P)$ turns each gap comparison into the
corresponding displayed integer inequality.
\end{proof}

Lemma~\ref{lem:chain} permits a reachability test rather than explicit
enumeration of maximal chains.  Start from the minimal vertices satisfying
the first inequality, traverse exactly the cover edges satisfying the middle
inequality, and accept a bad chain when a reached maximal vertex satisfies the
last inequality.  A presentation qualifies for the strict condition in the
definition of $m_2(4)$ precisely when this search finds no path.  Notice that
equality with $4$ is treated correctly: it makes the strict source condition
fail.

All computations for the theorem use unsigned $64$-bit integers.  For
$n\leq19$, extension counts are at most $n!$; the nonnegative terms in each
rank sum partition the extensions according to the insertion state, so a
rank sum is at most $n\,n!$.  The endpoint quantity is at most
$(n+1)n!$.  These bounds are below $2^{64}$ throughout the implemented range,
so neither the dynamic programs nor the comparisons overflow.

\section{Exhaustive computation}

The generator recursively emits every coordinatewise nondecreasing sequence
of pairs $(L_i,U_i)$ in the triangle $0\leq L_i\leq U_i\leq b$.  For every
$n=4,\ldots,17$ and every $1\leq a\leq\lfloor n/2\rfloor$, the program
applies Lemma~\ref{lem:encoding}, the dynamic programs
\eqref{eq:rankA}--\eqref{eq:rankB}, and Lemma~\ref{lem:chain}.  Table
\ref{tab:counts} reports the complete closed runs.  The implementation uses
NumPy~\cite{harris2020numpy} and Numba~\cite{lam2015numba}; these packages
accelerate the loops but are not mathematical or floating-point oracles.

\begin{table}[ht]
\centering
\caption{Complete two-chain-presentation counts.  The last column is zero at
every order.}\label{tab:counts}
\small
\begin{tabular}{r@{\qquad}r@{\qquad}r@{\qquad}r}
\toprule
$n$ & all presentations & width-two presentations & qualifiers\\
\midrule
4  & 30          & 20          & 0\\
5  & 65          & 50          & 0\\
6  & 301         & 260         & 0\\
7  & 714         & 651         & 0\\
8  & 3,312       & 3,150       & 0\\
9  & 8,397       & 8,142       & 0\\
10 & 39,094      & 38,457      & 0\\
11 & 104,005     & 102,982     & 0\\
12 & 484,705     & 482,196     & 0\\
13 & 1,337,219   & 1,333,124   & 0\\
14 & 6,227,729   & 6,217,822   & 0\\
15 & 17,678,834  & 17,662,451  & 0\\
16 & 82,204,044  & 82,164,842  & 0\\
17 & 238,819,349 & 238,753,814 & 0\\
\midrule
total & 346,907,798 & 346,767,961 & 0\\
\bottomrule
\end{tabular}
\end{table}

There is also a simple independent check on the width filter.  For fixed
$(a,b)$, the discarded width-one presentations have $L_i=U_i$ and are exactly
the weakly increasing length-$a$ sequences in $\{0,\ldots,b\}$, of which
there are $\binom{a+b}{a}=\binom{n}{a}$.  Thus the difference between the
second and third columns at order $n$ is
\[
 \sum_{a=1}^{\lfloor n/2\rfloor}\binom{n}{a},
\]
as the table data satisfy.  A separate dynamic count over the triangular
pair states also reproduces every entry in the presentation column.

Orders $1$--$3$ require no computation.  Order $1$ has width one.  At order
$2$ the unique width-two type is the antichain, whose singleton chains have
gap $3/2$.  At order $3$ the three width-two types are a two-element chain
plus an isolated point, the three-element V-poset, and its dual; each has a
maximal chain of gap at most $2$.  Combining these cases with the zero counts
in Table~\ref{tab:counts} proves Theorem~\ref{thm:main}.

\section{Independent verification and reproducibility}\label{sec:verification}

The primary program enumerates presentations, rather than canonical
isomorphism classes.  A second verifier deliberately changes all three
decisive components for $n\leq6$:
\begin{enumerate}
\item it enumerates subsets of forward relations, retains the transitive
relations of width exactly two, and canonicalizes over all permutations;
\item it recomputes extension counts and all rank sums by dynamic programming
over arbitrary order ideals; and
\item it explicitly enumerates every minimal-to-maximal cover path and tests
its gaps.
\end{enumerate}
The canonical domains from this verifier and from the two-chain generator are
identical.  Their numbers of unlabeled width-two posets for
$n=2,3,4,5,6$ are respectively
\[
  1,\ 3,\ 9,\ 25,\ 74.
\]
Every extension count, every rank sum, and every chain decision agrees on all
presentations through order $6$.  A separate audit also reran both full closed
computations, independently checked the presentation counts, and found the
same zero-qualifier totals.

The ancillary bundle contains the following decisive files:
\begin{itemize}
\item \texttt{evidence/enumerate\_width2.py}, the primary generator and exact
test;
\item \texttt{evidence/search-through-16.json} and
\texttt{evidence/search-order-17.json}, the closed run records;
\item \texttt{evidence/verify\_width2.py}, the independent small-order
verifier; and
\item \texttt{evidence/independent-audit-through-6.json}, its closed record.
\end{itemize}
From the bundle root, the computations can be replayed with
\begin{verbatim}
python3 evidence/enumerate_width2.py 16 > /tmp/search-through-16.json
python3 evidence/enumerate_width2.py 17 --min-n 17 > /tmp/search-order-17.json
(cd evidence && python3 verify_width2.py > /tmp/independent-through-6.json)
\end{verbatim}
The output JSON records exact integer counts; elapsed-time fields are
machine-dependent and play no role in the theorem.

\section{Limitations and next computation}

The proof establishes only $m_2(4)\geq18$.  The direct presentation generator
has $1{,}108{,}202{,}512$ presentations at order $18$, so extending the same
computation is materially larger.  No order-$18$ qualifier or other witness
is supplied here.  Moreover, presentation duplication is harmless for the
lower bound but would not meet the original question's requirement to give a
canonical, complete list of all minimizers.  A full resolution therefore
still requires an order-$18$ search (or a stronger structural argument), a
qualifying construction at the first possible order, and canonical
deduplication of every equality-order witness.

\section*{Computational and AI assistance disclosure}

OpenAI Codex assisted with formalizing the finite question, developing and
auditing the enumeration and verification programs, conducting the bounded
literature search, and preparing an initial manuscript.  Language-model
output is not used as mathematical evidence.  The theorem rests on the
explicit reductions above, exact integer computations, stored run records,
and independently reconstructing checks.  No proof assistant was used.

\section*{Data and code availability}

All code and finite run records needed for the reported result are included
in the ancillary bundle under the paths listed in
Section~\ref{sec:verification}.  No external dataset is used.

\bibliographystyle{plain}
\bibliography{references}

\end{document}
